% =============================================================================
% APENDICE A: ESPECIFICACION COMPLETA DE VARIABLES
% (fusion del antiguo apendice-lista con la descripcion detallada, julio 2026)
% =============================================================================

\section{Especificaci\'on Completa de Variables}
\label{sec:appendix_variables}

Este ap\'endice documenta el conjunto completo de 92 variables base empleadas en el an\'alisis, todas obtenidas de Bloomberg Terminal. A continuaci\'on se describe cada una de las siete categor\'ias que las organizan, su justificaci\'on te\'orica y las variables espec\'ificas que las componen.

\paragraph{\'Indices y ETFs de mercado (18 variables).}
La din\'amica de retornos entre clases de activos, regiones geogr\'aficas y sectores econ\'omicos puede aportar informaci\'on predictiva relevante. Los flujos de capital entre mercados desarrollados y emergentes, la rotaci\'on sectorial y los movimientos cambiarios reflejan cambios en el apetito por riesgo que frecuentemente anticipan los retornos del mercado amplio. Esta categor\'ia incluye (i)~\'indices de referencia norteamericanos, a saber, SPY (S\&P 500), QQQ (Nasdaq-100) e IWM (Russell 2000), que capturan la dispersi\'on entre \hyperlink{glos:large_small_caps}{\textit{large caps}}, tecnolog\'ia y \textit{small caps}; (ii)~ETFs internacionales, a saber, EFA (mercados desarrollados ex-US), EEM (mercados emergentes), VGK (Europa), EWJ (Jap\'on) y FXI (China), que reflejan la din\'amica global de riesgo; (iii)~ETFs sectoriales del S\&P 500, a saber, XLF (financiero), XLK (tecnolog\'ia), XLE (energ\'ia), XLV (salud), XLI (industrial) y XLU (servicios p\'ublicos), que capturan la rotaci\'on sectorial; y (iv)~divisas, a saber, DXY (\'indice d\'olar), EUR/USD y USD/JPY, que reflejan las condiciones de liquidez global y los diferenciales de pol\'itica monetaria. Se incluye adem\'as el \'indice FTSE NAREIT (bienes ra\'ices), sensible a las condiciones crediticias.

\paragraph{Indicadores macroecon\'omicos (19 variables).}
Las condiciones macroecon\'omicas fundamentales ejercen una influencia directa sobre los flujos de caja corporativos esperados y, por tanto, sobre la valoraci\'on de las acciones. \citet{chen1986} demostraron que variables como la producci\'on industrial, la inflaci\'on y los \hyperlink{glos:spread}{\textit{spreads}} de cr\'edito explican una porci\'on significativa de la variaci\'on transversal de retornos accionarios, estableciendo las bases de los modelos multifactoriales macroecon\'omicos. Esta categor\'ia comprende indicadores de actividad real (PIB trimestral QoQ y YoY, producci\'on industrial, utilizaci\'on de capacidad y ventas minoristas); indicadores del mercado laboral (n\'ominas no agr\'icolas o \textit{Non-Farm Payrolls}, tasa de desempleo, solicitudes iniciales de subsidio de desempleo e ingreso personal); indicadores de precios (IPC interanual y \hyperlink{glos:core_pce}{\textit{Core PCE}}); indicadores de confianza (sentimiento del consumidor de U. Michigan y confianza del consumidor del Conference Board); indicadores del sector vivienda (inicios de construcci\'on o \textit{Housing Starts}, ventas de vivienda nueva y existente); y los \'indices PMI manufacturero y de servicios (ISM), junto con el \'indice L\'ider Compuesto del Conference Board.

\paragraph{Tasas de inter\'es y curva de rendimiento (20 variables).}
La estructura temporal de tasas de inter\'es es uno de los predictores macroecon\'omicos m\'as robustos documentados en la literatura. \citet{harvey1989} demostr\'o que la pendiente de la curva de rendimiento anticipa recesiones econ\'omicas, mientras que \citet{fama1989} establecieron que los \textit{spreads} de cr\'edito y de t\'ermino predicen retornos accionarios a horizontes intermedios. Esta categor\'ia abarca la tasa de pol\'itica monetaria (Fed Funds); la curva de rendimiento completa del Tesoro estadounidense, que incluye \hyperlink{glos:treasury}{\textit{T-Bills}} a 3, 6 y 12 meses, y bonos a 2, 5, 10 y 30 a\~nos; los \'indices de rendimiento gen\'erico de Bloomberg para los vencimientos a 10 y 2 a\~nos (USGG10YR y USGG2YR), que a diferencia de los bonos individuales \hyperlink{glos:on_the_run}{\textit{on-the-run}} reflejan un compuesto de emisiones cercanas a cada plazo; los \textit{spreads} de t\'ermino precalculados (10Y--2Y y 10Y--3M); la tasa interbancaria a corto plazo (Libor 3M) y la \hyperlink{glos:swap}{tasa \textit{swap}} a 10 a\~nos; los \'indices de \hyperlink{glos:cds}{\textit{credit default swaps}}, a saber, CDX \textit{Investment Grade} y CDX \hyperlink{glos:ig_hy}{\textit{High Yield}}, que capturan la percepci\'on de riesgo crediticio corporativo; y los \'indices de bonos Bloomberg (US Aggregate Total Return, US Aggregate OAS, US Corporate Total Return), junto con el ETF LQD (\textit{investment grade} corporativo).

\paragraph{\textit{Commodities} (11 variables).}
Los precios de las materias primas operan como indicadores adelantados de presiones inflacionarias y cambios en la actividad econ\'omica global. Los movimientos en energ\'ia anticipan costos de producci\'on; los metales preciosos reflejan la demanda por activos refugio; y la actividad en los mercados de futuros de \textit{commodities} contiene informaci\'on predictiva sobre condiciones macroecon\'omicas y precios de activos \citep{hong2012}. Se incluyen futuros de energ\'ia (WTI o \textit{crude oil}, y \textit{heating oil}); metales preciosos tanto en futuros (oro, plata) como en ETFs (GLD, SLV); el ETF de petr\'oleo USO; \'indices amplios de \textit{commodities} (DBC, Bloomberg BCOM y S\&P GSCI); y el ETF DBA (agricultura).

\paragraph{Volatilidad (13 variables).}
La volatilidad impl\'icita constituye una medida directa de la incertidumbre percibida por el mercado y exhibe una relaci\'on contempor\'anea negativa con los retornos accionarios ampliamente documentada en la literatura, el denominado \textit{leverage effect}. \citet{ang2006} muestran, adem\'as, que la exposici\'on al riesgo de volatilidad agregada se refleja en la secci\'on cruzada de retornos esperados. La estructura temporal de la volatilidad, en particular la relaci\'on entre VIX de corto y largo plazo, captura cambios en el r\'egimen de riesgo que frecuentemente preceden movimientos direccionales del mercado. Esta categor\'ia comprende el VIX y sus variantes por horizonte temporal, a saber, VIX1M (1 mes), VIX3M (3 meses) y VXV (3 meses, variante), cuya estructura temporal permite identificar inversiones indicativas de estr\'es; \'indices de volatilidad impl\'icita para otras clases de activos (VXEEM para mercados emergentes, VXEFA para mercados desarrollados, GVZ para oro, OVX para petr\'oleo, TYVIX para bonos del Tesoro y MOVE para el mercado de renta fija); el CVIX (divisas); el futuro VIX del mes pr\'oximo (VIY1); y el V2X (Euro Stoxx 50), que captura el r\'egimen de volatilidad europeo.

\paragraph{Sentimiento de mercado (10 variables).}
El sentimiento inversor, definido como la propensi\'on a especular no justificada por los fundamentales, genera desviaciones sistem\'aticas de precios respecto a su valor intr\'inseco. \citet{baker2006} demostraron que per\'iodos de alto sentimiento preceden retornos bajos, especialmente en acciones dif\'iciles de arbitrar, lo cual ofrece poder predictivo en la secci\'on cruzada de retornos. Esta categor\'ia incluye la encuesta AAII (\textit{American Association of Individual Investors}) de sentimiento alcista; ratios \hyperlink{glos:put_call}{\textit{put/call}} de opciones sobre acciones (CBOE Equity Put/Call y Put Index); indicadores de \hyperlink{glos:breadth}{\textit{breadth}} de mercado (NYSE \hyperlink{glos:new_highs_lows}{\textit{New Highs--New Lows}}, \hyperlink{glos:advance_decline}{\textit{Advance-Decline Line}}, TICK y TRIN o \hyperlink{glos:tick_trin}{\textit{Arms Index}}); el oscilador McClellan y su \'indice de sumaci\'on, que capturan la amplitud del mercado en m\'ultiples horizontes; y el \'indice SRVOL (\textit{Short-Range Volatility}), que captura la volatilidad de corto plazo.

\paragraph{Datos OHLCV del SPY y tasa libre de riesgo.}
Finalmente, se incluyen los cinco campos est\'andar del registro de precio diario del ETF SPY, a saber, apertura, m\'aximo, m\'inimo, cierre y volumen (\hyperlink{glos:ohlcv}{OHLCV}), que constituyen la materia prima para el c\'alculo de indicadores t\'ecnicos, estimadores de volatilidad OHLC y caracter\'isticas de microestructura. La tasa libre de riesgo, derivada del \textit{T-Bill} a 3 meses, se utiliza para el c\'alculo de retornos en exceso que conforman la variable objetivo del sistema.
